\textbf{Phần 1. Trắc nghiệm lựa chọn}
\setcounterpageref{ex}{0}

\begin{ex} %[3D2V54]
Tính giá trị của phép tính: $4528 + 3617$.
\choice
{$8135$}
{\True $8145$}
{$8155$}
{$7145$}
\loigiai{
Đặt tính và cộng lần lượt từ phải sang trái: $4528 + 3617 = 8145$.
\textbf{Đáp án của câu hỏi là B}
}
\end{ex}

\begin{ex} %[3D2V54]
Một trường tiểu học có $1250$ học sinh nam và $1340$ học sinh nữ. Trường đó có tất cả bao nhiêu học sinh?
\choice
{$2580$ học sinh}
{$2600$ học sinh}
{$2490$ học sinh}
{\True $2590$ học sinh}
\loigiai{
Tổng số học sinh: $1250 + 1340 = 2590$ học sinh.
\textbf{Đáp án của câu hỏi là D}
}
\end{ex}

\begin{ex} %[3D2V54]
Tìm số tự nhiên $x$, biết: $x - 3450 = 2860$.
\choice
{\True $6310$}
{$6210$}
{$6410$}
{$5310$}
\loigiai{
$x = 2860 + 3450 = 6310$.
\textbf{Đáp án của câu hỏi là A}
}
\end{ex}

\begin{ex} %[3D2V54]
Tính nhanh giá trị biểu thức: $1250 + 3400 + 750 + 600$.
\choice
{$5800$}
{$6200$}
{\True $6000$}
{$5900$}
\loigiai{
Nhóm thuận tiện: $(1250 + 750) + (3400 + 600) = 2000 + 4000 = 6000$.
\textbf{Đáp án của câu hỏi là C}
}
\end{ex}

\begin{ex} %[3D2V54]
Một đội làm đường ngày thứ nhất làm được $1450\text{ m}$, ngày thứ hai làm được nhiều hơn ngày thứ nhất $350\text{ m}$. Cả hai ngày đội đó làm được bao nhiêu mét đường?
\choice
{$1800\text{ m}$}
{\True $3250\text{ m}$}
{$3150\text{ m}$}
{$3300\text{ m}$}
\loigiai{
Ngày hai làm: $1450 + 350 = 1800\text{ m}$. Cả hai ngày: $1450 + 1800 = 3250\text{ (m)}$.
\textbf{Đáp án của câu hỏi là B}
}
\end{ex}

\begin{ex} %[3D2V54]
Một kho thóc có $3250\text{ kg}$ thóc. Người ta nhập thêm hai đợt thóc: đợt một nhập $1800\text{ kg}$, đợt hai nhập $2450\text{ kg}$. Trong kho có tất cả bao nhiêu ki-lô-gam thóc?
\choice
{$6500\text{ kg}$}
{$7400\text{ kg}$}
{$7600\text{ kg}$}
{\True $7500\text{ kg}$}
\loigiai{
Tổng số thóc: $3250 + 1800 + 2450 = 7500\text{ (kg)}$.
\textbf{Đáp án của câu hỏi là D}
}
\end{ex}

\begin{ex} %[3D2V54]
Nếu một số hạng tăng thêm 250 đơn vị và số hạng kia tăng thêm 450 đơn vị thì tổng sẽ
\choice
{\True Tăng thêm 700 đơn vị}
{Tăng thêm 200 đơn vị}
{Tăng thêm 600 đơn vị}
{Không đổi}
\loigiai{
Tổng mới tăng thêm: $250 + 450 = 700$ đơn vị.
\textbf{Đáp án của câu hỏi là A}
}
\end{ex}

\begin{ex} %[3D2V54]
Tìm tổng của số lớn nhất có 4 chữ số khác nhau và số bé nhất có 4 chữ số khác nhau.
\choice
{$10\,900$}
{$10\,889$}
{\True $10\,899$}
{$11\,000$}
\loigiai{
Số lớn nhất: 9876. Số bé nhất: 1023. Tổng: $9876 + 1023 = 10\,899$.
\textbf{Đáp án của câu hỏi là C}
}
\end{ex}

\begin{ex} %[3D2V54]
Một trang trại có 2560 con gà và số vịt nhiều hơn số gà là 480 con. Trang trại có tất cả bao nhiêu con gà và vịt?
\choice
{$3040$ con}
{\True $5600$ con}
{$5500$ con}
{$5620$ con}
\loigiai{
Số vịt: $2560 + 480 = 3040$ con. Tổng cộng: $2560 + 3040 = 5600$ con.
\textbf{Đáp án của câu hỏi là B}
}
\end{ex}

\begin{ex} %[3D2V54]
Điền chữ số thích hợp vào dấu $*$ để: $4*56 + 2341 = 6897$. Chữ số $*$ là
\choice
{$4$}
{$6$}
{$3$}
{\True $5$}
\loigiai{
Ta có: $4*56 = 6897 - 2341 = 4556 \Rightarrow * = 5$.
\textbf{Đáp án của câu hỏi là D}
}
\end{ex}

\begin{ex} %[3D2V54]
Bố mua một chiếc quạt điện giá $450\,000$ đồng và một ấm đun nước giá $280\,000$ đồng. Hỏi bố phải trả tất cả bao nhiêu tiền?
\choice
{\True $730\,000$ đồng}
{$720\,000$ đồng}
{$740\,000$ đồng}
{$750\,000$ đồng}
\loigiai{
Tổng số tiền: $450\,000 + 280\,000 = 730\,000$ đồng.
\textbf{Đáp án của câu hỏi là A}
}
\end{ex}

\begin{ex} %[3D2V54]
Một thư viện trường học tuần thứ nhất cho mượn 1420 cuốn sách, tuần thứ hai cho mượn nhiều hơn tuần thứ nhất 160 cuốn sách. Cả hai tuần thư viện cho mượn bao nhiêu cuốn sách?
\choice
{$2840$ cuốn sách}
{$3100$ cuốn sách}
{\True $3000$ cuốn sách}
{$2900$ cuốn sách}
\loigiai{
Tuần hai cho mượn: $1420 + 160 = 1580$ cuốn. Cả hai tuần: $1420 + 1580 = 3000$ cuốn sách.
\textbf{Đáp án của câu hỏi là C}
}
\end{ex}

\textbf{Phần 2. Trắc nghiệm đúng sai}
\setcounterpageref{ex}{0}

\begin{ex} %[3D2V54]
Xét tính đúng sai của các khẳng định về phép cộng trong phạm vi 10 000:
\choiceTF
{\True Phép cộng các số có 4 chữ số thực hiện lần lượt từ phải sang trái (từ hàng đơn vị).}
{Tổng của hai số có 4 chữ số luôn luôn là một số có 4 chữ số.}
{\True Khi đổi chỗ các số hạng trong một tổng thì tổng không thay đổi.}
{$3500 + 4500 = 7000$.}
\loigiai{
\begin{itemize}
    \item Ý a) Đúng: Quy tắc đặt tính và tính.
    \item Ý b) Sai: Sai vì ví dụ $6000 + 5000 = 11\,000$ có 5 chữ số.
    \item Ý c) Đúng: Tính chất giao hoán của phép cộng.
    \item Ý d) Sai: Sai vì $3500 + 4500 = 8000$.
\end{itemize}
\begin{itemize}
    \item a) \textbf{Đúng}.
    \item b) \textbf{Sai}.
    \item c) \textbf{Đúng}.
    \item d) \textbf{Sai}.
\end{itemize}
}
\end{ex}

\begin{ex} %[3D2V54]
Một xã có hai thôn: thôn Đông có 3420 nhân khẩu, thôn Đoài có 2850 nhân khẩu:
\choiceTF
{\True Cả xã có tất cả 6270 nhân khẩu.}
{\True Thôn Đông nhiều hơn thôn Đoài 570 nhân khẩu.}
{Làm tròn số dân thôn Đông đến hàng trăm ta được 3400 người.}
{\True Tổng số dân làm tròn đến hàng trăm là khoảng 6300 người.}
\loigiai{
\begin{itemize}
    \item Ý a) Đúng: $3420 + 2850 = 6270$ người.
    \item Ý b) Đúng: $3420 - 2850 = 570$ người.
    \item Ý c) Sai: Sai vì hàng chục là 2 nên đúng là 3400 người (mệnh đề đúng: làm tròn thành 3400). Sửa câu sai: Làm tròn thôn Đông đến hàng trăm được 3500 người.
    \item Ý d) Đúng: $6270$ làm tròn đến hàng trăm là 6300.
\end{itemize}
\begin{itemize}
    \item a) \textbf{Đúng}.
    \item b) \textbf{Đúng}.
    \item c) \textbf{Sai}.
    \item d) \textbf{Đúng}.
\end{itemize}
}
\end{ex}

\textbf{Phần 3. Trả lời ngắn}
\setcounterpageref{ex}{0}

\begin{ex} %[3D2V54]
Tính kết quả phép cộng: $4835 + 2748$.
\shortans{7583}
\loigiai{
\textbf{Đáp số:} $7583$\\
$4835 + 2748 = 7583$.
}
\end{ex}

\begin{ex} %[3D2V54]
Tìm $x$, biết: $x - 2480 = 3520$.
\shortans{6000}
\loigiai{
\textbf{Đáp số:} $6000$\\
$x = 3520 + 2480 = 6000$.
}
\end{ex}

\begin{ex} %[3D2V54]
Tính nhanh giá trị biểu thức: $2400 + 1750 + 600 + 250$.
\shortans{5000}
\loigiai{
\textbf{Đáp số:} $5000$\\
$(2400 + 600) + (1750 + 250) = 3000 + 2000 = 5000$.
}
\end{ex}

\begin{ex} %[3D2V54]
Một cửa hàng buổi sáng bán được $1350\text{ kg}$ gạo, buổi chiều bán được nhiều hơn buổi sáng $420\text{ kg}$ gạo. Hỏi cả hai buổi cửa hàng bán được bao nhiêu ki-lô-gam gạo?
\shortans{3120}
\loigiai{
\textbf{Đáp số:} $3120$\\
Buổi chiều: $1770\text{ kg}$. Cả hai buổi: $1350 + 1770 = 3120\text{ (kg)}$.
}
\end{ex}

\begin{ex} %[3D2V54]
Tính tổng của số chẵn lớn nhất có 3 chữ số và số tròn nghìn bé nhất có 4 chữ số.
\shortans{1998}
\loigiai{
\textbf{Đáp số:} $1998$\\
$1000 + 998 = 1998$.
}
\end{ex}

\begin{ex} %[3D2V54]
Một vườn ươm có 3500 cây keo và 2800 cây bạch đàn. Hỏi vườn ươm có tất cả bao nhiêu cây giống?
\shortans{6300}
\loigiai{
\textbf{Đáp số:} $6300$\\
$3500 + 2800 = 6300$ cây.
}
\end{ex}
